\section{Выводы}

В этой лабораторной я впервые всерьёз профилировал свою программу и понял, что представление
о том, где она тратит время, часто неверное. Я ожидал, что дороже всего будут повороты в
AVL-дереве, а оказалось, что почти половину времени съедал {\ttfamily istringstream}, который я
поставил ради удобства разбора строки. Из этого вывод: перед оптимизацией надо измерять, а не
угадывать.

{\ttfamily gprof} удобен и прост, но видит только мой код. Сумма по профилю (0.35~с) почти вдвое
меньше реального времени (0.89~с), и без {\ttfamily callgrind} я бы не узнал, куда делась остальная
половина. Поэтому {\ttfamily gprof} лучше сочетать с инструментом, который считает и библиотечный код.

{\ttfamily dmalloc} и {\ttfamily memcheck} утечек в программе не нашли, зато с их помощью я увидел
\emph{число} выделений памяти: 213 тыс. на 50 тыс. слов. Отсюда пришла идея хранить ключ прямо в
узле одной аллокацией. {\ttfamily dmalloc} оказался менее удобен для C++: для {\ttfamily new} он не
показывает место в исходнике, только для {\ttfamily malloc} после подключения {\ttfamily dmalloc.h}.
Зато он работает без отдельной виртуальной машины и в режиме {\ttfamily high} проверяет всю кучу.
{\ttfamily massif} нагляднее всего показал, что {\ttfamily Load} временно удваивает память.

Самые важные найденные ошибки были не в структуре данных, а в разборе ввода: программа принимала
неверные команды, сохраняла файлы, которые сама не могла загрузить, и молча очищала словарь на
опечатку в команде. Тесты из условия этого не ловили, нашёл я это только целенаправленными
граничными тестами. Ещё понял, что запись файла поверх старого надо делать через временный файл
и {\ttfamily rename}, иначе ошибка посередине записи уничтожает данные.

В итоге исправленная версия работает в 1.7--2 раза быстрее, тратит на 33--41\% меньше памяти, а
{\ttfamily Load} ускорился в 7.6 раза. Эти навыки пригодятся в любой программе, где важны скорость
и память: сначала профилировщик и проверка памяти, потом исправления, потом повторный замер.
\pagebreak
